\section{Q4: Constructing Class Groups}

In question 1, we have established that if $p \nmid h(\Qd)$, then coprime solution $(x, y)$ for the equation 
\begin{equation}\label{eq:4main}
	y^p = x^2 + r^2 d
\end{equation}
with $r^2 d \equiv 2 \mod 4$ satifies $y \geq d$.
Therefore, by constructing coprime solutions to \eqref{eq:4main} such that $y < d$, we can show that $p | h(\Qd)$.

Let us consider two equations 
\begin{align}
	3^p &= 5^2 + k_1 \label{contra1}\\
	3^p &= 7^2 + k_2 \label{contra2}
\end{align}
where $k_1 = 3^p - 5^2$ and $k_2 = 3^p - 7^2$.
We have the following congruence relations
\begin{align}
	3^a &\equiv 3 \mod 4   &\text{for any odd }a\\
	5^2 &\equiv 1 \mod 4 &\\
	7^2 &\equiv 1 \mod 4 &
\end{align}
therefore $k_1 \equiv k_2 \equiv 2 \mod 4$ and \eqref{contra1} and \eqref{contra2} satisfies our contrains for question 1. 

Write $k_1 = r_1^2 d_1$ where $d_1$ are square-free.
Since $ k_1 \equiv 2 \mod 4$, we have $2 | k_1$ while $4 \nmid k_1$, so $2 |d_1$. 
If $\frac{k_1}{2}$ is not a perfect square, there is another prime $q \neq 2$ that divides $d_1$, so $d_1 > 3$, and our solution $(x = 5, y = 3)$ for \eqref{contra1} satisfies $y < d_1$ which means $p | h(\Q(\sqrt{-d_1}))$.
The same arguments apply for $k_2$ and \eqref{contra2} as well.

Note that $\frac{k_1}{2} + 12 = \frac{k_2}{2}$ and we have 
\begin{align}
	7^2 - 6^2 &= 13 \\ 
	(a+1)^2 - a^2 &> 13 \text{ for any } a \geq 6
\end{align}
Therefore, if $\frac{k_1}{2} > 36 = 6^2$, it is impossible that both of $\frac{k_1}{2}$ and $\frac{k_2}{2}$ are perfect squares, and our $p$ divides at least one of $h(\Q(\sqrt{-d_1}))$, $h(\Q(\sqrt{-d_2}))$

It is clear that $k_1 > 36$ for all $p \geq 5$. 
So our arguments above works for all $p \geq 5$.
The case $p=3$ leaves no problem as $h(\Q(\sqrt{-26})) = 6$ is divisible by 3.
In this way we have shown that for any odd prime $p$, there exists some $d$ such that $p | h(\Qd)$.

Moreover, in either case we have $k_i < 3 ^p $. As $d| k_i$
it implies that 
$$
d < 3^p.
$$
Taking logarithm base 3 on both sides gives us
$$
p > \frac{\log(d)}{\log(3)} 
$$
We know $p | h(\Qd)$, meaning $h(\Qd) \geq p$. Combining the above two inequalities gives us
$$
h(\Qd) > \frac{\log(d)}{\log(3)}
$$
as desired.
